\documentclass[conference]{IEEEtran}
\IEEEoverridecommandlockouts
\usepackage{cite}
\usepackage{amsmath,amssymb}
\usepackage{booktabs}
\usepackage{url}
\usepackage{array}
\usepackage[hidelinks]{hyperref}
\hypersetup{pdftitle={Next-Technique Prediction in an APT29 Evaluation Sequence: A Pilot Benchmark},pdfauthor={}}

\title{Next-Technique Prediction in an APT29 Evaluation Sequence: A Pilot Benchmark}
\author{\IEEEauthorblockN{Anonymous Submission}}

\begin{document}
\maketitle

\begin{abstract}
Forecasting an adversary's next action during an intrusion is a useful goal for threat-informed defense, but empirical evaluation is constrained by the scarcity of public, ordered attack traces. This paper presents a preliminary benchmark of next-technique prediction on one ordered sequence reconstructed from the MITRE ATT\&CK Evaluations APT29 material. We cross-check the reconstructed sequence across five vendor result files and find identical ordered substep/technique pairs, while retaining the limitation that these files represent the same emulation flow rather than independent campaigns. The historical experiment record describes a global-frequency baseline, a hierarchical co-occurrence baseline, first- and second-order Markov models, and two LSTM training regimes; the original training code and per-position prediction logs are unavailable, so these historical scores are not independently reproducible from the materials retained for this revision. A synthetic corpus of 500 tactic-constrained sequences is used for training, with a chronological split of 98 observed APT29 steps and 42 held-out steps. The historical table reports 4/42 top-3 hits for the mixed-training LSTM and 7/42 top-5 hits for the first-order Markov model. These values are retained only as unverified historical figures pending rerun of a fully specified implementation; they must not be interpreted as independently reproduced results. These results are descriptive rather than evidence of reliable generalization: the evaluation contains one campaign-derived sequence, few test positions, dependent sequential predictions, and a synthetic generator that uses the target technique inventory. We characterize the study as a pilot benchmark and specify the validation needed before operational claims are justified.
\end{abstract}

\begin{IEEEkeywords}
MITRE ATT\&CK, adversarial technique prediction, sequence modeling, Markov model, LSTM, cybersecurity evaluation, synthetic data
\end{IEEEkeywords}

\section{Introduction}
MITRE ATT\&CK provides a shared vocabulary for describing adversary tactics and techniques and is widely used in threat intelligence, detection engineering, and incident response \cite{roy2023sok}. A more ambitious task is to predict which technique may follow a partially observed sequence. Such predictions could help analysts prioritize investigation or prepare detections, but should not be treated as deterministic forecasts: ATT\&CK tactics express adversary goals, and actual operations can branch, repeat, omit, or reorder behaviors \cite{mitreTactics}.

Prior ATT\&CK research has studied technique associations and co-occurrence across reported groups and malware, including clustering and association-rule analyses \cite{alshaer2020,rahman2026}. These approaches can expose relationships between techniques, but association is not equivalent to temporal predictability. Research has also examined the broader use of ATT\&CK, telemetry-to-technique mapping, and automated extraction from threat-intelligence reports \cite{jiang2025,lupinacci2026,fayyazi2024,haque2026}. The availability of execution-ordered, public ground truth remains a practical constraint for directly evaluating next-step prediction.

This paper narrows the question to a controlled, preliminary comparison of frequency-based, association-based, transition-based, and neural sequence models on one public APT29 evaluation sequence. It makes three limited contributions:
\begin{itemize}
\item It defines next-technique prediction as a ranked top-$k$ task over a fixed 58-technique vocabulary.
\item It documents a proposed synthetic-to-real evaluation protocol and records a historical comparison of six method configurations. Because the original implementation and raw predictions are unavailable, the historical comparison is treated as provisional rather than reproduced evidence.
\item It documents the limitations that prevent the reported numbers from establishing cross-campaign or operational performance.
\end{itemize}

The study does not claim a deployable predictor or cross-campaign generalization. Its purpose is to make the data and validation gaps concrete and to provide a cautious starting point for a more reproducible benchmark.

\section{Related Work}
\subsection{ATT\&CK Association and Sequence Modeling}
Al-Shaer \emph{et al.} studied statistical associations among ATT\&CK techniques and explored clustering as a way to infer unobserved techniques from observed ones \cite{alshaer2020}. Erfan \emph{et al.} studied co-occurrence patterns using adversary and malware technique data \cite{rahman2026}. Both lines of work motivate association-based baselines, while leaving a distinct evaluation question: whether ordered prefixes provide useful information for predicting the next observed technique.

A systematization of ATT\&CK research identifies broad application areas and validation gaps \cite{roy2023sok}; a later review analyzes 417 peer-reviewed publications and discusses challenges in ATT\&CK's application to dynamic environments \cite{jiang2025}. Other work studies mapping low-level telemetry to ATT\&CK \cite{lupinacci2026} and using language models to interpret or extract tactics and techniques from reports \cite{fayyazi2024,haque2026}. These are related but different tasks: mapping observations to technique labels or extracting labels from reports does not itself evaluate one-step-ahead sequence forecasting.

Sequential threat forecasting also predates the recent technique-level work. Ghafir \emph{et al.} combine alert correlation with a hidden Markov model to infer and predict APT stages from security alerts \cite{ghafir2019}. Pappu studies smoothed higher-order Markov models for ATT\&CK tactic prediction, addressing sparse transition counts with discounting and multiple random seeds \cite{pappu2025}. These studies use different prediction units and data protocols from the present technique-level top-$k$ task; their reported metrics are therefore not directly comparable, but they establish relevant prior art.

Most directly, Raj \emph{et al.} introduce a hybrid LSTM--Markov framework for attack-chain forecasting and risk assessment, reporting 4,849 campaign chains and 8,437 operational intrusion flows, together with risk-ranked multi-step forecasts \cite{raj2026}. Their publisher notes that the severity analysis has 33 independent campaign-level labels, rather than treating all derived chains as independent severity outcomes. This contemporary work means that the present paper should not claim to be the first or only study of ATT\&CK sequence forecasting. The distinction here is narrower: a top-$k$ comparison on one ordered APT29 evaluation sequence with a 98/42 chronological split. The reported metrics are not directly comparable across the two studies because their corpora, split definitions, targets, and evaluation protocols differ. A head-to-head comparison on the same held-out sequence would be needed before claiming a model advantage over the hybrid approach.

\subsection{Synthetic Attack Data and Emulation Sources}
APTGen generates artificial targeted-attack sequences from existing incident reports and executes corresponding activity to create logs and labels \cite{takahashi2020}. The present study uses a simpler tactic-constrained random-walk generator and does not reproduce APTGen's report-grounded sequence construction or execution-based log generation. The real sequence is reconstructed from public APT29 evaluation results \cite{mitreApt29,shen2024}. That reconstruction is a single evaluation artifact, not a sample of independent campaigns.

\section{Task and Evaluation Questions}
Let $\mathcal{T}=\{t_1,\ldots,t_V\}$ be the fixed technique vocabulary, with $V=58$. Let the reconstructed ordered sequence be $S=(s_1,\ldots,s_n)$, where $s_i\in\mathcal{T}$. The held-out suffix is $Y=(y_1,\ldots,y_m)=(s_{n-m+1},\ldots,s_n)$, with $m=42$. At test position $j$, the model receives the observed prefix through $s_{n-m+j-1}$ and returns a ranked list $R_j^{(k)}$ of $k$ distinct candidate techniques. Top-$k$ accuracy is
\begin{equation}
 A_k=\frac{1}{m}\sum_{j=1}^{m}\mathbb{1}\{y_j\in R_j^{(k)}\}.
\end{equation}

The pilot study addresses three descriptive questions: (RQ1) do order-aware methods achieve higher top-$k$ accuracy than order-agnostic baselines on this split? (RQ2) does greater model complexity improve the observed ranking metrics? (RQ3) how much can be concluded when synthetic training data and a single real sequence are used? These questions are evaluated on one split only; no claim of population-level inference is made.

\section{Methodology}
\subsection{Real Sequence and Vocabulary}
The reconstruction uses the public MITRE ATT\&CK Evaluations APT29 results, flattening the structured technique-step records and sorting substeps by (major step number, branch letter, numeric suborder) to obtain 140 technique instances drawn from a 58-technique vocabulary \cite{mitreApt29,mitreJoystick}. The original evaluation material describes a representative operational flow; the reconstructed order should therefore be understood as an emulation/evaluation sequence, not a complete or uniquely canonical chronology for all APT29 campaigns.

We performed a cross-vendor reconstruction check using the McAfee, Microsoft, FireEye, CrowdStrike, and Bitdefender result JSON files in MITRE's public ATT\&CK evaluation data repository \cite{mitreJoystick}. For each file, flattening the structured steps and applying the sorting rule above produced 140 unique substep labels. The complete ordered sequence of (substep label, technique ID) pairs matched exactly across all five files. This check supports the reproducibility of the ordering operation and reduces concern that the result is an artifact of one vendor export. It does not create five independent samples: the files report the same MITRE-designed emulation flow. The model evaluation therefore remains a single-sequence pilot and cannot establish cross-campaign generalization.

\subsection{Synthetic Corpus}
The historical manuscript describes a corpus of 500 generated sequences, lengths sampled between 20 and 140 steps, a realized mean length of 75.5, and 37,274 next-technique training instances after window decomposition. These three statistics have not been reconciled against the original generator, which is unavailable. They are therefore reported as historical descriptions, not verified corpus properties; a new run must regenerate and report its own exact counts. The experiment-specific tactic order is Execution $\rightarrow$ Persistence $\rightarrow$ Privilege Escalation $\rightarrow$ Defense Evasion $\rightarrow$ Credential Access $\rightarrow$ Discovery $\rightarrow$ Lateral Movement $\rightarrow$ Collection $\rightarrow$ Command and Control $\rightarrow$ Exfiltration $\rightarrow$ Impact. A random walk advances one tactic with probability 0.55, repeats the current tactic with probability 0.25, moves backward one or two positions with probability 0.13, or skips forward two or three positions with probability 0.07. A technique is then sampled from the pool associated with the selected tactic; the reported in-scope pool contains no Impact technique, so Impact does not yield sampled instances.

This order is a modeling heuristic, not an ATT\&CK-mandated sequence. ATT\&CK defines tactics by the adversary's goal or ``why'' for using a technique; it does not require all techniques to follow one fixed linear order \cite{mitreTactics}. The generator also uses technique-to-tactic membership and the inventory in scope for the target evaluation. Although the ordered held-out suffix is not described as being used to fit transition probabilities, the synthetic distribution is not independent of the target vocabulary. This is a form of structural target-domain information and limits the interpretation of the transfer experiment.

\subsection{Models}
The historical description lists (1) global frequency, which ranks techniques by training frequency, and (2) hierarchical co-occurrence clustering, reportedly built from a $\pm5$-position co-occurrence window, a count-to-distance transformation, and average-linkage clustering into 12 clusters. The clustering ranking rule is not fully specified and cannot be recovered from the retained artifacts; accordingly, its historical top-$k$ scores are not reproducible as currently documented. The transition models are a first-order Markov chain and a second-order Markov chain with backoff to first order under sparsity; both use additive smoothing. The LSTM uses 32-dimensional embeddings, 64 hidden units, dropout 0.2, and a six-technique input window. For the clustering method, the precise rule that converts cluster membership into a fully ranked list of candidates must be specified, because top-$3$ and top-$5$ metrics depend on that ranking.

Two historical LSTM training regimes are described. In mixed training, synthetic examples and the first 98 real steps were reportedly combined for 60 epochs. In pretrain-then-fine-tune, the LSTM was reportedly pretrained on synthetic data for 40 epochs and fine-tuned on the first 98 real steps for 60 epochs at a reduced learning rate. These settings are not verified against source code and must not be assumed to describe the new clean-room implementation. A reproducible artifact should specify the random seeds, optimizer and learning-rate values, smoothing constants, software versions, and the full ranking and tie-breaking rules used to derive top-$3$ and top-$5$ candidates from every baseline. The present description does not define all of these details.

\subsection{Chronological Holdout}
The reported split uses the first 98 of 140 steps as observed context and training/fine-tuning data and the final 42 steps for evaluation. At each test position, the model receives the true preceding prefix, including earlier held-out steps after they have occurred. This is a teacher-forced, one-step-ahead evaluation: it tests prediction when the most recent observed technique is available, not a free-running forecast of a complete future sequence. It avoids using later ground-truth steps to make an earlier prediction, but test predictions are sequentially dependent and should not be treated as 42 independent experiments.

\section{Reported Results}
Table~\ref{tab:results} reports the stated top-$k$ scores alongside the corresponding hit counts implied by the rounded accuracies over 42 test positions. The counts are arithmetic conversions and should be checked against per-position model outputs in a reproducible release.

\begin{table}[t]
\centering
\caption{Reported top-$k$ accuracy on the 42-step APT29 holdout. Parenthetical values give the corresponding hit counts out of 42.}
\label{tab:results}
\scriptsize
\setlength{\tabcolsep}{3.3pt}
\begin{tabular}{p{0.45\columnwidth}ccc}
\toprule
Method & Top-1 & Top-3 & Top-5\\
\midrule
Global frequency & 0.024 (1/42) & 0.048 (2/42) & 0.071 (3/42)\\
Co-occurrence clustering & 0.024 (1/42) & 0.048 (2/42) & 0.071 (3/42)\\
Markov, 2nd-order & 0.000 (0/42) & 0.071 (3/42) & 0.143 (6/42)\\
Markov, 1st-order & 0.000 (0/42) & 0.071 (3/42) & 0.167 (7/42)\\
LSTM, pretrain + fine-tune & 0.024 (1/42) & 0.071 (3/42) & 0.095 (4/42)\\
LSTM, mixed training & 0.000 (0/42) & 0.095 (4/42) & 0.119 (5/42)\\
\bottomrule
\end{tabular}
\end{table}

Three descriptive observations follow from the historical table only. First, the displayed order-aware models have higher top-3 and top-5 scores than the frequency and clustering rows. Second, the mixed-training LSTM has the best top-3 score (4/42, approximately 9.5\%), while the first-order Markov model has the best top-5 score (7/42, approximately 16.7\%). No reported model correctly ranks the next technique first more than once; the first-order Markov and mixed-training LSTM rows have zero top-1 hits. Third, the absolute results remain modest. Under uniform random ranking over 58 distinct techniques, the reference hit rates would be $1/58\approx1.7\%$ for top-1, $3/58\approx5.2\%$ for top-3, and $5/58\approx8.6\%$ for top-5. These are reference values, not a substitute for a frequency-matched or transition-matched baseline.

Because only aggregate metrics are available in the manuscript, we cannot compute paired uncertainty intervals or a valid significance test from the table alone. A naive independent-binomial interval would also understate the dependence induced by sequential evaluation. Raw per-position predictions and repeated runs are needed for uncertainty analysis. The observed top-$k$ differences should therefore be treated as descriptive rather than statistically established superiority.

\section{Discussion}
The historical table is consistent with the hypothesis that predecessor context may carry signal beyond global technique prevalence, but because the scores have not been reproduced, this interpretation is tentative and does not establish that sequential models generalize across campaigns. The first-order Markov model's top-5 result is stronger than that of the reported LSTM configurations on this split, while the mixed-training LSTM's top-3 score is higher. Neither pattern is evidence that one model family is generally superior: the number of test positions is small, only one sequence is used, and the reported table does not include variation across random seeds or campaigns.

A second implication is methodological. Synthetic sequence volume is not equivalent to independent real-world evidence. A generator that samples techniques uniformly within tactic pools can impose the desired tactic progression while failing to capture real transition frequencies, branching, campaign-specific behavior, or observation noise. Moreover, ATT\&CK tactic membership is a structural prior derived in part from the target evaluation inventory. A stronger design should compare at least three settings: (i) technique pools derived without target-sequence information; (ii) pools based on an independently versioned ATT\&CK knowledge base; and (iii) target-informed pools, clearly labeled as such. This would quantify the effect of structural target-domain information rather than treating synthetic-to-real transfer as a single condition.

For practical forecasting, a free-running evaluation should also be reported separately from teacher-forced one-step prediction. In teacher forcing, every test prediction receives the correct preceding test steps. In a rollout, later inputs include the model's own predictions, so errors can compound. The former is a useful local metric; it should not be presented as equivalent to forecasting an unseen campaign trajectory.

\section{Operationalization of the Benchmark}
\subsection{Ranking Protocol and Baseline Interpretation}
A top-$k$ metric is meaningful only when every model returns a complete, consistently ordered candidate list. Frequency baselines rank techniques by descending training count, with deterministic tie-breaking required to make reruns stable. A first-order Markov model estimates the next technique conditional on the immediately preceding technique; a second-order model conditions on the last two observed techniques and backs off when the relevant context is unseen or too sparse. Additive smoothing prevents zero-probability transitions but does not resolve the underlying sparsity caused by a small sequence and a 58-label vocabulary. For neural models, the ranking is induced by the output scores over the same vocabulary. Comparisons are fair only if all methods use the same candidate inventory, the same observed prefix at each prediction point, and the same exclusion rules for already-seen or invalid labels.

The retained historical clustering baseline requires particular caution. Its description mentions a local co-occurrence window, a distance transformation, and average-linkage clustering, but it does not specify how clusters are converted into ranked candidates at inference time. Several plausible rules could produce different top-$k$ lists even with identical clusters. For that reason, the historical clustering row cannot be reconstructed from its method label alone. The accompanying clean-room implementation uses a documented association-ranking baseline instead; this is a related but distinct method and must be named accordingly in any updated result table. Preserving this distinction prevents an implementation substitution from being mistaken for an exact replication.

\subsection{Metrics and Reference Points}
Top-$k$ accuracy measures whether the observed next technique appears anywhere in the first $k$ predictions. It does not measure calibration, the quality of the full ranking, the cost of false alarms, or whether a prediction would have enabled a defender to intervene. Top-1 is the strictest of the three metrics reported here, while top-3 and top-5 reward broader candidate sets. Since the vocabulary has 58 techniques, increasing $k$ mechanically increases the chance of including the target. Consequently, top-$k$ values should be interpreted together, not selectively reported at the cutoff that favors a particular model.

The uniform-ranking reference rates of $k/58$ are useful only as a simple scale check. They are not a strong operational baseline because techniques are not equally prevalent and sequence transitions are not uniform. A more informative benchmark should include global frequency, previous-technique persistence, first-order transition frequency, and a frequency-matched random ranker. For each seed, the evaluation should retain the exact candidate ranking at every position, not only an aggregate score. Where several sequences are available, scores should be reported per sequence and macro-averaged across campaigns so that a long trace does not dominate the summary. With only one sequence, a narrow percentage-point difference should not be presented as a stable estimate of expected performance.

\subsection{Reproducible Execution Record}
A benchmark run should record the software environment, model configuration, data-source revision, all random seeds, and hashes of input files. It should also export the generated sequences, the ordered real sequence, each test prefix, the true next technique, each model's top-five candidates, and the aggregate counts used to compute the metrics. The manifest should distinguish immutable source material from generated artifacts. Failed downloads, missing labels, empty technique pools, and training failures should be surfaced as explicit errors rather than silently replaced with fallback data. These records make it possible to audit not only the final score but also whether two methods were evaluated on exactly the same examples.

For a multi-seed synthetic experiment, the seed is a nuisance variable that can affect the generated training corpus and optimization path. Reporting only the best seed would introduce selection bias. The planned rerun should report every prespecified seed and summarize mean and dispersion, while retaining per-seed outputs. Seed variation does not substitute for campaign variation: three training runs on the same APT29 sequence remain three runs of one evaluation target, not three independent campaign observations. A later benchmark should prioritize independent traces and a held-out-campaign design over merely increasing the number of synthetic sequences.

\subsection{From Offline Ranking to Defensive Use}
The current task predicts one technique label from an already observed sequence. Operational use would require additional steps not measured here: mapping alerts or telemetry to reliable technique labels, handling missing and duplicated observations, representing parallel or branching activity, estimating uncertainty under distribution shift, and deciding how a ranked candidate list should affect analyst workflow. An incorrect prediction could waste investigation time or create false confidence, especially when the model assigns a high score to a plausible but unsupported technique. Any deployment-oriented evaluation should therefore measure alert burden, analyst utility, and robustness to incomplete or noisy observations, alongside predictive metrics.

A useful next-stage design would compare teacher-forced next-step prediction with a separate free-running rollout. In the teacher-forced setting used by this pilot, the true preceding held-out technique becomes available for the next prediction. In a rollout, the model's earlier predictions become part of later input, allowing errors to compound. These protocols answer different questions and should be reported separately. A further ablation should remove target-informed technique pools from synthetic generation, or derive pools from an independently versioned knowledge base, to quantify how much apparent transfer depends on structural information about the evaluation target. None of these additional evaluations has been completed in the present draft; they are specified as requirements for subsequent validation rather than results claimed here.

\section{Limitations and Required Validation}
The study has five material limitations. First, the evaluation contains one reconstructed APT29 emulation sequence. Its ordered step/technique pairs were identical across five vendor result files, but these are multiple exports of the same flow, not independent campaigns, so cross-campaign generalization remains untested. Second, the holdout comprises 42 sequentially dependent positions. Third, the target technique inventory informs synthetic tactic pools, so the training corpus is structurally target-informed. Fourth, the full ranking rule for the clustering baseline and several implementation details are not specified in the method description. Fifth, a reproducible release needs to include source code, raw per-position predictions, random seeds, the generated corpus or its generator, environment details, and the exact reconstructed sequence.

Before the results can support a stronger empirical claim, the implementation should be checked and released where permitted; the exact 140-step reconstructed sequence should be documented with its source and sorting rule; all candidate-ranking rules should be specified; counts in Table~\ref{tab:results} should be validated against per-step output; and the experiment should be repeated across multiple random seeds. At minimum, report per-position predictions, top-$k$ metrics, macro-averaged results if multiple sequences become available, calibration if probabilities are used, and uncertainty estimates that respect sequence dependence (for example, a block bootstrap only when enough sequence material supports it). The strongest improvement would be evaluation on multiple independently sourced ordered sequences, with leave-one-campaign-out or campaign-held-out testing. These checks are prerequisites for claims beyond this pilot.

\section{Reproducibility Status and Data Availability}
The public source sequence can be reconstructed from MITRE's structured evaluation result files using the source URL and sorting rule described above. The historical model implementation, generated synthetic corpus, model checkpoints, and per-position prediction rankings were not available when this revision was prepared. As a result, the historical scores and generator statistics are not independently verified. The accompanying clean-room code defines a new baseline implementation and exports its source hashes, generated sequences, ranked predictions, and metrics. It is not an exact recreation of the lost historical implementation: in particular, its co-occurrence association baseline is not equivalent to the historical hierarchical clustering method. Any results generated by that code must be presented as new experiments and substituted consistently throughout the manuscript. The paper should not be submitted as a validated empirical comparison until this rerun is completed and the reported results are updated.

\section{Conclusion}
This paper defines a cautious pilot protocol for next-technique ranking on one ordered MITRE ATT\&CK APT29 evaluation sequence. A cross-vendor reconstruction check found identical ordered substep/technique pairs across five public result files, strengthening confidence in the reconstruction procedure without broadening the sample beyond one emulation flow. The historical results table is retained as an unverified record because the original training code and per-position rankings are unavailable. It therefore cannot support claims of model superiority or operational utility. A submission making empirical claims must replace that table with metrics produced by the included clean-room implementation, report its exact generator and ranking rules, and evaluate additional independent sequences before making generalization claims.

\IEEEtriggeratref{8}
\begin{thebibliography}{15}
\bibitem{roy2023sok} S. Roy, E. Panaousis, C. Noakes, A. Laszka, S. Panda, and G. Loukas, ``SoK: The MITRE ATT\&CK Framework in Research and Practice,'' arXiv:2304.07411, 2023. [Online]. Available: \url{https://arxiv.org/abs/2304.07411}
\bibitem{mitreTactics} MITRE, ``Enterprise Tactics,'' MITRE ATT\&CK. [Online]. Available: \url{https://attack.mitre.org/tactics/}. Accessed: Oct. 11, 2026.
\bibitem{alshaer2020} R. Al-Shaer, J. M. Spring, and E. Christou, ``Learning the Associations of MITRE ATT\&CK Adversarial Techniques,'' in \emph{Proc. IEEE Conf. Communications and Network Security (CNS)}, 2020, pp. 1--9, doi: 10.1109/CNS48642.2020.9162207. [Online]. Available: \url{https://arxiv.org/abs/2005.01654}
\bibitem{rahman2026} M. Erfan, A. Ryan, L. Williams, and M. R. Rahman, ``Investigating Co-occurrences of MITRE ATT\&CK Techniques,'' \emph{Computers \& Security}, vol. 170, Art. no. 105042, 2026, doi: 10.1016/j.cose.2026.105042.
\bibitem{jiang2025} Y. Jiang, Q. Meng, F. Shang, N. Oo, L. T. H. Minh, H. W. Lim, and B. Sikdar, ``MITRE ATT\&CK Applications in Cybersecurity and The Way Forward,'' arXiv:2502.10825, 2025. [Online]. Available: \url{https://arxiv.org/abs/2502.10825}
\bibitem{lupinacci2026} M. Lupinacci, L. Arena, F. Blefari, and A. Furfaro, ``A Graph-Based Approach for Mapping Kernel-Level Telemetry to MITRE ATT\&CK,'' arXiv:2609.12841, 2026. [Online]. Available: \url{https://arxiv.org/abs/2609.12841}
\bibitem{fayyazi2024} R. Fayyazi, R. Taghdimi, and S. J. Yang, ``Advancing TTP Analysis: Harnessing the Power of Large Language Models with Retrieval Augmented Generation,'' arXiv:2401.00280, 2024. [Online]. Available: \url{https://arxiv.org/abs/2401.00280}
\bibitem{haque2026} M. N. Haque, S. Hamer, B. Wroblewski, M. R. Rahman, and L. Williams, ``Beyond Single Reports: Evaluating Automated ATT\&CK Technique Extraction in Multi-Report Campaign Settings,'' arXiv:2604.07470, 2026. [Online]. Available: \url{https://arxiv.org/abs/2604.07470}
\bibitem{ghafir2019} I. Ghafir, K. G. Kyriakopoulos, S. Lambotharan, F. J. Aparicio-Navarro, B. Assadhan, H. Binsalleeh, and D. M. Diab, ``Hidden Markov Models and Alert Correlations for the Prediction of Advanced Persistent Threats,'' \emph{IEEE Access}, vol. 7, pp. 99508--99520, 2019, doi: 10.1109/ACCESS.2019.2930200.
\bibitem{pappu2025} K. Pappu, ``Smoothed Markov Chains for MITRE ATT\&CK Prediction: Addressing Data Scarcity in Cyber Kill Chains,'' in \emph{Proc. 5th Int. Conf. on Electronic Information Engineering and Computer Communication (EIECC)}, 2025, pp. 761--764, doi: 10.1109/EIECC67963.2025.11409611.
\bibitem{raj2026} M. Raj, N. D. Bastian, L. Fiondella, and G. Kul, ``MITRE ATT\&CK-based Attack Chain Prediction Using Hybrid LSTM-Markov Models for Cyber Risk Assessment,'' in \emph{Proc. 23rd Int. Conf. on Security and Cryptography (SECRYPT)}, 2026, pp. 1039--1050, doi: 10.5220/0015075400004103.
\bibitem{takahashi2020} Y. Takahashi, S. Shima, R. Tanabe, and K. Yoshioka, ``APTGen: An Approach towards Generating Practical Dataset Labelled with Targeted Attack Sequences,'' in \emph{Proc. 13th USENIX Workshop on Cyber Security Experimentation and Test (CSET)}, 2020. [Online]. Available: \url{https://www.usenix.org/conference/cset20/presentation/takahashi}
\bibitem{mitreApt29} MITRE Engenuity, ``ATT\&CK Evaluations: APT29,'' 2019. [Online]. Available: \url{https://attackevals.mitre-engenuity.org/enterprise/APT29/}
\bibitem{shen2024} X. Shen, Z. Li, G. Burleigh, L. Wang, and Y. Chen, ``Decoding the MITRE Engenuity ATT\&CK Enterprise Evaluation: An Analysis of EDR Performance in Real-World Environments,'' in \emph{Proc. ACM Asia Conf. on Computer and Communications Security (ASIA CCS '24)}, 2024, pp. 96--111, doi: 10.1145/3634737.3645012.
\bibitem{mitreJoystick} MITRE, ``ATT\&CK Evaluations structured result data, APT29 (McAfee, Microsoft, FireEye, CrowdStrike, and Bitdefender JSON files),'' \emph{mitre-attack/joystick} public repository, archived Sep. 13, 2023. [Online]. Available: \url{https://github.com/mitre-attack/joystick/tree/8bee0d3128abc08c4e00c60b01c1608e162bb9ce/data/evaluations/apt29}. Accessed: Oct. 11, 2026.
\end{thebibliography}

\end{document}
